\documentclass[10pt,a4paper]{amsart}
\usepackage[margin=19mm]{geometry}
\usepackage{amsmath,amssymb,mathtools}
\usepackage[scr=rsfs,cal=euler]{mathalfa}
\usepackage{xfrac}
\usepackage[unicode,hidelinks,bookmarks=false]{hyperref}
\newcommand{\ie}{i.e.\ }
\newcommand{\cf}{cf.\ }
\newcommand{\resp}{respectively\ }
\newcommand{\Subsection}{\subsection}
\providecommand{\nobreakdash}{\nobreak\hbox{-}}
\setlength{\emergencystretch}{2em}
\multlinegap0pt
\usepackage{bm}
\usepackage{bbm}
\usepackage{dsfont}
\usepackage{xspace}
\usepackage{cancel}
\usepackage{mathtools}
\usepackage{amsthm}
\makeatletter
\@ifundefined{theorem}{\newtheorem{theorem}{Theorem}}{}
\@ifundefined{lemma}{\newtheorem{lemma}[theorem]{Lemma}}{}
\@ifundefined{proposition}{\newtheorem{proposition}[theorem]{Proposition}}{}
\makeatother
\title[Quantum Mechanics on Non-Hausdorff One-Manifolds with Finite]{Quantum Mechanics on Non-Hausdorff One-Manifolds with Finite Graph Resolutions: Analytic Completion, Branching, and Invariant-Sector Scattering}
\author{Abhiram Sripat}
\date{Source: arXiv:2607.26080v3 (CC BY 4.0)}
\begin{document}
\maketitle

\noindent\emph{Source and licence.} Quantum Mechanics on Non-Hausdorff One-Manifolds with Finite Graph Resolutions: Analytic Completion, Branching, and Invariant-Sector Scattering. Version arXiv:2607.26080v3 (CC BY 4.0), distributed under the Creative Commons Attribution 4.0 International licence, as recorded in the arXiv licence metadata for this exact version and shown on the abstract page https://arxiv.org/abs/2607.26080v3. Full rights notice: licenses/arxiv-2607.26080-CC-BY-4.0.txt.\par\smallskip

\noindent\textit{Editorial scope.} Counted unit: the Proposition whose source label is \texttt{prop:quantitative-cutoff-flatness} in the article (source0.tex lines 13404--13441, source label \texttt{prop:quantitative-cutoff-flatness}), delivered together with the article's own complete original proof of it (source0.tex lines 13443--13502, one \texttt{proof} environment of 60 lines). No mathematics is written, repaired or re-derived: statement and proof are the article's own text line for line. Required context retained uncounted: lem:angular-cutoff-derivative-bounds (lines 13277--13306); lem:normal-taylor-decay (lines 13342--13368). Every definition, display and label of those lines is retained unchanged. Omitted from this package and not redistributed: 1--13276, 13307--13341, 13369--13403, 13442--13442, 13503--13879. Nothing inside a retained excerpt is omitted or altered: every retained body is delivered exactly as the article's lines are written, so no omission receipt of any kind accompanies this package. Citations: the retained excerpts carry no citation command at all, so no bibliography entry of the article is transcribed and no reference is redistributed. Reference adaptation: none was needed and none was made; every reference inside the delivered unit resolves inside it, so no unresolved reference or citation is produced. Printed slips kept literal: no printed word, symbol, hypothesis or number of the counted statement or of its proof was altered. Layout and class: the article's own class is \texttt{article} with packages including fontenc, inputenc, lmodern, microtype, amsmath, amssymb, amsfonts, amsthm, amscd, mathtools, mathrsfs, geometry, titling, enumitem; the wrapper is the lane's pinned \texttt{amsart} editorial wrapper at 10pt on A4, which restates the theorem-like environments the retained text needs and adds the article's own macro definitions that the retained text needs (none), with extra wrapper packages (bm, bbm, dsfont, xspace, cancel, mathtools). The delivered numbering is the unit's own. Assets: no retained span contains a figure environment, a graphics-inclusion command, a PDF-page inclusion, a table or a TikZ picture; no image file is copied into this package, the package declares no asset dependency and no image file is redistributed with it. Licence: version arXiv:2607.26080v3 of this article is distributed under the Creative Commons Attribution 4.0 International licence, as recorded in the arXiv licence metadata for this exact version and shown on the abstract page https://arxiv.org/abs/2607.26080v3. A full rights notice is delivered at licenses/arxiv-2607.26080-CC-BY-4.0.txt. Characters: every retained excerpt is pure ASCII, so the delivered engine, XeLaTeX with its default Latin Modern text font, reports no missing character.
\begin{lemma}[Derivative bounds for angular cutoffs]
\label{lem:angular-cutoff-derivative-bounds}
Let
\[
  K\Subset S.
\]
For every pair of multi-indices \(\alpha,\beta\), there exists
\[
  C_{K,\alpha,\beta}>0
\]
such that
\[
  \boxed{
  \left|
    \partial_p^\alpha
    \partial_v^\beta
    \chi_a(p,v)
  \right|
  \leq
  C_{K,\alpha,\beta}
  |v|^{-|\beta|}
  }
\]
whenever
\[
  p\in K,
  \qquad
  0<|v|<2\eta(p).
\]
\end{lemma}

\begin{lemma}[Normal Taylor decay]
\label{lem:normal-taylor-decay}
Let
\[
  r\in\Gamma^k(\mathcal U,F)
\]
and suppose that its full ambient \(k\)-jet vanishes along \(S\):
\[
  j_S^k r=0.
\]
Then, in every adapted coordinate chart and local bundle frame,
\[
  \boxed{
  \partial_p^\alpha
  \partial_v^\beta r(p,v)
  =
  o
  \left(
    |v|^{k-|\alpha|-|\beta|}
  \right)
  }
\]
locally uniformly in \(p\), whenever
\[
  |\alpha|+|\beta|\leq k.
\]
\end{lemma}

\begin{proposition}[Quantitative cutoff-flatness]
\label{prop:quantitative-cutoff-flatness}
For every pair of multi-indices \(\alpha,\beta\) satisfying
\[
  |\alpha|+|\beta|\leq k,
\]
one has
\[
  \boxed{
  \partial_p^\alpha
  \partial_v^\beta u_a(p,v)
  =
  o
  \left(
    |v|^{k-|\alpha|-|\beta|}
  \right)
  }
\]
locally uniformly in \(p\).

The zero extension
\[
  [\chi_ar]_0(x)
  :=
  \begin{cases}
    \chi_a(x)r(x),&x\notin S,\\
    0,&x\in S
  \end{cases}
\]
belongs to
\[
  \Gamma^k(\mathcal U,F)
\]
and satisfies
\[
  j_S^k[\chi_ar]_0=0.
\]
\end{proposition}

\begin{proof}
Leibniz' rule gives
\[
  \partial_p^\alpha
  \partial_v^\beta
  (\chi_ar)
  =
  \sum_{
    \substack{
      \alpha_1+\alpha_2=\alpha\\
      \beta_1+\beta_2=\beta
    }
  }
  C_{\alpha_1,\alpha_2,\beta_1,\beta_2}
  \left(
    \partial_p^{\alpha_1}
    \partial_v^{\beta_1}\chi_a
  \right)
  \left(
    \partial_p^{\alpha_2}
    \partial_v^{\beta_2}r
  \right).
\]
By
Lemma~\ref{lem:angular-cutoff-derivative-bounds},
the first factor is
\(O \left( |v|^{-|\beta_1|} \right)\).
By
Lemma~\ref{lem:normal-taylor-decay},
the second factor is
\(o \left( |v|^{ k-|\alpha_2|-|\beta_2| } \right)\).
Their product is therefore
\[
  o
  \left(
    |v|^{
      k-|\alpha_2|-|\beta_1|-|\beta_2|
    }
  \right)
  =
  o
  \left(
    |v|^{
      k-|\alpha_2|-|\beta|
    }
  \right).
\]
Since
\(|\alpha_2|\leq|\alpha|\),
this implies
\(o \left( |v|^{k-|\alpha|-|\beta|} \right)\).
Summing the finitely many Leibniz terms proves the estimate.

In particular, every derivative through order \(k\) tends
continuously to zero along the zero section. Defining all those
derivatives to be zero on \(S\), and inducting on the derivative order
in adapted coordinates, shows that the zero extension is \(C^k\).
Every derivative through order \(k\) vanishes on \(S\), so its
\(k\)-jet there is zero.
\end{proof}

\end{document}
